\chapter{Differential forms and Stokes}

\paragraph{Notations} Some common notations in this chapter are, if \(U
\subseteq \mathbb{C} \) is an open set, \(f \in C^{1}{(U)}\), then
\[
  \frac{\partial f}{\partial z} := \frac{1}{2}{\left( \frac{\partial f}{\partial
  x} - i \frac{\partial f}{\partial y} \right)} \quad ; \quad \frac{\partial
f}{\partial \bar{z}} := \frac{1}{2} {\left( \frac{\partial f}{\partial x} +
i \frac{\partial f}{\partial y}\right)} \quad ; \quad df := \frac{\partial f}{\partial z} dz + \frac{\partial f}{\partial
\bar{z}} d\bar{z}\] 

If \(\varphi : V \to U\) is holomorphic, then
\[
  \frac{\partial f \circ \varphi}{\partial z} = \varphi'{(z)} {\left( \frac{\partial
  f}{\partial z} \circ \varphi \right)}  \quad ; \quad \frac{\partial f \circ
\varphi}{\partial \bar{z}} = \overline{\varphi'{(z)}} {\left( \frac{\partial
f}{\partial \bar{z}} \circ \varphi \right)} 
\]
\begin{definition}{}
    A \(1\)-form on \(U\) is \(\omega = p{(z)} \,dz + q {(z)}
    \,d\bar{z}\) and
    \[
      \varphi^{*} \omega \overset{\text{def}}{=} p{(\varphi {(z)})} \varphi'{(z)} dz +
      q{(\varphi{(z)})} \overline{\varphi'{(z)}} d\bar{z}
    \]
\end{definition}

\begin{remark}{}
    \[
      d{(f \circ \varphi)} = \varphi^{*}{(df)}
    \]
\end{remark}
\begin{definition}{}
    Let \({(U_{\alpha} , \varphi _{\alpha} )}\) be an atlas on \(X\). Then
\begin{itemize}[label = \bullet]
    \item A 1-form \(\omega\) is a collection \(\omega_{\alpha} \) of 1-form on
        \(\varphi _{\alpha} {(U_{\alpha} )} \subseteq \mathbb{C} \) 
    \item \(\omega_{\alpha_{2}} = {\left( \varphi _{\alpha_{1}} \circ \varphi
        _{\alpha_{2}}^{-1} \right)}^{*} \omega_{\alpha_{1}}  \) whenever
        \(U_{\alpha_{1}} \cap U_{\alpha_{2}} \neq \varnothing\) 
\end{itemize}
    Such that \(\omega\) is well defined on \(X\).

\end{definition}


\begin{definition}{2-forms}
    If \(\eta = f{(z)}dz \wedge d\bar{z}\) (on \(U \subseteq \mathbb{C} \)) and
    \(\varphi : V \to U\) is \emph{holomorphic}, then
    \[
      \varphi ^{*} \eta = f \circ \varphi \left| \varphi '{(z)} \right| ^2 dz
      \wedge d\bar{z}
    \]
\end{definition}
\begin{remark}{}
    Let \(\omega = p dz + q d\bar{z}\) be a 1-form. Then
    \[
      d {\left( \varphi ^{*} \omega \right)} = {\left( \frac{\partial \varphi
      }{\partial z} \circ q  - \frac{\partial p}{\partial \bar{z}}
\circ \varphi \right)} \left| \varphi ' \right| ^2 dz \wedge d\bar{z}
\implies d {\left( \varphi ^{*}\omega \right)} = \varphi ^{*}{(d\omega)}
    \]
\end{remark}


\begin{definition}{2-forms}
    Let \({(U_{\alpha} , \varphi _{\alpha} )}\) be an atlas on \(X\). Then
\begin{itemize}[label = \bullet]
    \item A 2-form \(\eta\) is a collection \(\eta_{\alpha} \) of 2-forms on
        \(\varphi _{\alpha} {(U_{\alpha} )} \subseteq \mathbb{C} \) 
    \item \(\eta_{\alpha_{2}} = {\left( \varphi _{\alpha_{1}} \circ \varphi
        _{\alpha_{2}}^{-1} \right)}^{*} \eta_{\alpha_{1}}  \) whenever
        \(U_{\alpha_{1}} \cap U_{\alpha_{2}} \neq \varnothing\) 
\end{itemize}
    Such that \(\eta\) is well defined on \(X\).
\end{definition}

We will denote by
\begin{itemize}[label = --]
    \item \(\Omega^{0}{(X)}\) the set of \(C^{\infty}\) functions on
\(X\)
    \item \(\Omega^{1}{(X)}\) the set of \(C^{\infty}\) 1-forms on \(X\) 
    \item \(\Omega^{2}{(X)}\) the set of \(C^{\infty}\) 2-forms on \(X\) 
\end{itemize}

The operator \(d\) is well defined
\[
  \Omega^{0}{(X)} \overset{d}{\to} \Omega^{1}{(X)} \overset{d}{\to}
  \Omega^{2}{(X)}
\]
and moreover \(d\circ d = 0\) 

\paragraph{Formulas:} if I take \(f \in \Omega^{0}{(X)}\), then I can set
locally \(\partial f = \frac{\partial f}{\partial z} dz\) and
\(\bar{\partial} f = \frac{\partial f}{\partial \bar{z}}
d\bar{z}\). Similarly, if \(\omega = p\, dz + q\, d\bar{z}\) is a 1-form,
then I can set locally
\[
  \partial \omega = \frac{\partial q}{\partial z} dz \wedge d\bar{z} \quad
  ; \quad \bar{\partial} \omega = -\frac{\partial p}{\partial \bar{z}}
  dz \wedge d\bar{z}
\]
So we have \(d = \partial + \bar{\partial}\) and \(d^2 = 0 = \partial^2 =
\bar{\partial}^2\).

For \emph{1-forms}, we have a decomposition. Let us denote by
\(\Omega^{{(1,0)}}{(X)}\) the set of 1-forms such that \(\omega = p{(z)}dz\)
locally and similarly by \(\Omega^{{(0,1)}}{(X)}\) the set of 1-forms such that
locally \(\omega = q{(z)} d\bar{z}\). If \(p\) is holomorphic (meromorphic), then \(\omega\)
is called holomorphic (meromorphic).

The following diagram holds:
\[
\begin{tikzcd}
	& {\Omega^0(X)} & \\
	{\Omega^{(1,0)}(X)} && {\Omega^{(0,1)}(X)} \\
	{\{0\}} & {\Omega^2(X)} & {\{0\}}
	\arrow["\partial"', from=1-2, to=2-1]
	\arrow["{\bar{\partial}}", from=1-2, to=2-3]
	\arrow["\partial"', from=2-1, to=3-1]
	\arrow["{\bar{\partial}}", from=2-1, to=3-2]
	\arrow["\partial"', from=2-3, to=3-2]
	\arrow["{\bar{\partial}}", from=2-3, to=3-3]
\end{tikzcd}
\]

\section{Integration and differential forms}

Let \(\omega \in \Omega^{1}{(X)}\) and \(C : [0, 1] \to X\) be smooth. We want
to get to a sensible definition of \(\int _C \omega\).

Pick a partition \(0 = t_{0} < t_{1} < \dots < t_{n} = 1\) of \([0, 1]\) such
that \(\forall k = 0, \dots, n-1\) we have \(C{([t_k, t_{k+1}])} \subseteq
{(U_k, \varphi _k)} \) a chart. At this point we can define
\begin{equation}\label{eq:defin_int}
    \int _C \omega = \sum_{k=0}^{n-1} \int_{\underbrace{\varphi_k {(C[t_k,
    t_{k+1} ])}}_{\text{curve in \(\mathbb{C}\)}}}
    \overbrace{{\left( \varphi ^{-1}_k \right)}^{*} \omega}^\text{1-form on
    \(\varphi _k{(U_k)} \subseteq \mathbb{C} \)} 
\end{equation}
which in particular \emph{does not} depend on \emph{charts} and
\emph{subdivision}. 

In \textbf{local coordinates}, we have
\[
  \int_{\gamma} {\left( p \,dx + \varphi \,dy \right)} \overset{\text{def}}{=}
  \int_{0}^{1} {\left( {\left( p \circ \gamma{(t)} \right)} \gamma_1'{(t)} +
  {\left( q \circ \gamma{(t)} \right)} \gamma'_2 {(t)} \right)} \, dt
\]
where \(\gamma = {(\gamma_{1}, \gamma_{2})}: [0, 1] \to \mathbb{R}^2\).

For 2-forms it's a bit harder.

Let \(K \subseteq X \) be \emph{compact}, \(\eta \in \Omega^2{(X)}\) with 
\({(U_{i}, \varphi_i)}_{i \in I} \) a locally finite atlas on \(X\). Let
\({(p_{i})}_{i \in I} \) be a \emph{partition of unity} such that
\begin{enumerate}[label = \arabic*., nosep]
    \item \(\forall i\), \(p_{i} \in C^{\infty}{(X)}\), 
        \(\mathrm{supp}{(p_{i})} \subseteq U_{i} \) and \(p_{i} \ge 0\) 
    \item \(\forall x \in X\), \(\sum_{i \in I} p _i {(x)} = 1\) 
\end{enumerate}
Then we can define
\[
  \int_K \eta = \sum_{i \in I} \int_{\varphi_i {(K \cap U_{i})}} {\left( p_{i}
  \circ \varphi_i ^{-1} \right)} {\left( \varphi_i ^{-1} \right)}^{*} \eta
\]
which does not depend on \({(p_{i})}_{i \in I} \) or \({(U_{i}, \varphi_i)}_{i
\in I}\).

Let \(K \subseteq X \) compact such that \(\mathring{K} \neq \varnothing\) and
\(\partial K\) is a finite disjoint union of simple loops, i.e. \(\partial K\)
is a 1-dimensional submanifold of \(X\).

\begin{proposition}{}
    Riemann Surfaces are orientable
\end{proposition}
Hence \(\partial K\) is naturally oriented.

\begin{theorem}{Stokes}
    Let \(K \subseteq X \) as in the previous paragraph and \(\omega \in
    \Omega^{1}{(X)}\), then
    \[
      \int_K d\omega = \int _{\partial K} \omega
    \]
\end{theorem}

\begin{corollary}[Consequences of Stokes]
    Let \(X\) compact, \(f \in \Omega^{0}{(X)}\) and \(\omega \in
    \Omega^{1}{(X)}\). Then
\begin{itemize}[label = --, nosep]
    \item if \(\displaystyle \omega \in \Omega^{{(0, 1)}{(X)}}\), then \(\displaystyle \int_X
        f \, \partial \omega = - \int _X \partial f \wedge \omega\) 
    \item if \(\displaystyle \omega \in \Omega^{{(1,0)}}{(X)}\), then
        \(\displaystyle \int_X f \, \bar{\partial} \omega  = - \int _X
        \bar{\partial} f \wedge \omega\) 
\end{itemize}
\end{corollary}

\begin{proof}{}
    In general, \(d{(f \omega)} = {(df)} \wedge \omega + f \, d\omega\).
    Integrating, we get
    \[
        \int _X f \, d\omega = \underbrace{\int_X d{(f \omega)}}_\text{= 0 by
        Stokes \((\partial X = \varnothing)\) }   - \int_X {(df)} \wedge \omega 
    \]
    and if \(\omega \in \Omega^{{(0, 1)}}{(X)}\), then \(d\omega =
    \partial\omega\) and (not obvious) \({(df)} \wedge \omega = \partial f
    \wedge \omega\).  % TODO: prof wrote f in WW !!
    Analogously with \(\omega \in \Omega^{{(1, 0)}{(X)}}\).

\end{proof}

\begin{definition}{Laplacian}
    Let \(f \in \Omega^{0}{(X)}\). Then the \textbf{laplacian}
    \[
      \Delta f := 2i \bar{\partial} \partial f = - 2i \partial \bar{\partial} f
      \in  \Omega^2{(X)}
    \]
    which in local charts is
    \[
      \Delta f = - {\left( \frac{\partial^2 f}{\partial x ^2} + \frac{\partial
      ^2 f}{\partial y^2}\right)} dx \wedge dy
    \]
\end{definition}
\begin{note}{}
    \(\Delta f\) is a 2-form, not a function
\end{note}


Which leads to this rephrasing of Stokes:
\begin{theorem}{Green's formula}
    For all \(f, g \in \Omega^{0}{(X)}\) (\(X\) compact), we have
    \[
      \int_X {(\Delta f)} g = \int_X f \Delta g = -\frac{1}{2i} \int_X
      \bar{\partial} f \wedge \partial g
    \]
\end{theorem}

\begin{proof}{}
    Integration by parts.
\end{proof}

\begin{corollary}{}
    Assume \(X\) compact connected. Let \(f \in \Omega^{0}{(X)}\) be such that
    \(\Delta f = 0\). Then \(f\) is constant.
\end{corollary}

\begin{proof}{}
    Integrating \({(\Delta f)} \bar{f}\) we get
    \[
      0 = \int_X {(\Delta f)} \bar{f} = - \frac{1}{2i} \int_X \bar{\partial} f
      \wedge \partial \bar{f} = \frac{1}{2i} \int_X \partial f \wedge
      \bar{\partial} \bar{f}
    \]
    and in local coordinates, \(\bar{\partial} f \wedge \partial \bar{f} =
    \left| \frac{\partial f}{\partial \bar{z}} \right| ^2 d\bar{z}\). Since
    \(\int_X {(\Delta f)} \bar{f}\) is a sum of positive terms, we get that
    locally \(\frac{\partial f}{\partial \bar{z}} = 0\). Hence, \(f\) is
    holomorphic, thus constant.
\end{proof}

